\section{Comandos úteis e boas práticas}\label{sec:boas-praticas}

\subsection{Links, endereços e datas}

Use \lstinline|\url{https://www.latex-project.org}| para exibir um endereço completo: \url{https://www.latex-project.org}. Para associar um texto ao endereço, use \lstinline|\href{https://www.latex-project.org}{Projeto LaTeX}|, que produz \href{https://www.latex-project.org}{Projeto LaTeX}. O comando \lstinline|\today| imprime a data atual segundo o idioma configurado: \today.

\subsection{Espaçamento e tamanho}

O \LaTeX{} calcula automaticamente a maior parte do espaçamento. Quando um ajuste pontual for indispensável, \lstinline|\vspace{5mm}| acrescenta espaço vertical e \lstinline|\hspace{5mm}| acrescenta espaço horizontal. Os comandos \lstinline|\small|, \lstinline|\footnotesize|, \lstinline|\large| e \lstinline|\Large| alteram o tamanho a partir do ponto em que aparecem; limite seu efeito com chaves, por exemplo, \lstinline|{\small texto menor}|.

\subsection{Criação de comandos}

Comandos repetitivos podem ser definidos uma vez e reutilizados. Este projeto define \lstinline|\fonteimagem| em \texttt{artigo-abnt.cls}:

\begin{lstlisting}[language={[LaTeX]TeX},caption={Definição e uso de um comando com argumento},label={cod:comando}]
\newcommand{\fonteimagem}[1]{%
  \par\smallskip{\footnotesize #1}\par
}

\fonteimagem{Fonte: elabora\c{c}\~ao pr\'opria.}
\end{lstlisting}

O número \lstinline|[1]| informa quantos argumentos o comando recebe, e \lstinline|#1| representa o primeiro argumento.

\subsection{Checklist de revisão}

Antes de entregar o documento:

\begin{enumerate}
  \item verifique se todo \lstinline|\begin| possui o \lstinline|\end| correspondente;
  \item mantenha chaves e colchetes balanceados;
  \item use rótulos únicos e sem espaços;
  \item confirme se todas as citações existem no arquivo \texttt{.bib};
  \item remova pacotes, figuras e arquivos que não são usados;
  \item faça uma compilação completa e revise avisos, sumário, listas e referências.
\end{enumerate}

Este documento é simultaneamente um guia e um exemplo compilável. Para adaptar o modelo, substitua os textos demonstrativos, mantenha apenas as seções necessárias e consulte o próprio código-fonte sempre que quiser reproduzir algum elemento.
